\documentclass[noauthor,nooutcomes,hints,handout]{ximera}

%% page layout
\usepackage[in,headings]{fullpage}
\raggedright
\setlength\headheight{13.6pt}


%% fonts
\usepackage{euler}

\usepackage{FiraMono}
\renewcommand\familydefault{\ttdefault} 
\usepackage{mathastext}
\usepackage[htt]{hyphenat}

\usepackage[T1]{fontenc}
\usepackage[scaled=1]{FiraSans}

\usepackage{wedn}
\usepackage[T1]{fontenc}

%% wrap text around scripts
\usepackage{wrapfig}

\tikzset{>=stealth}
%% snap! scripts
\usepackage{scratch3}

\usepackage{adjustbox}

%% journal style
\makeatletter
\newcommand\journalstyle{%
  \def\activitystyle{activity-chapter}
  \def\maketitle{%
    \addtocounter{titlenumber}{1}%
                {\flushleft\small\sffamily\bfseries\@pretitle\par\vspace{-1.5em}}%
                {\flushleft\LARGE\sffamily\bfseries\thetitlenumber\hspace{1em}\@title \par }%
                {\vskip .6em\noindent\textit\theabstract\setcounter{question}{0}\setcounter{sectiontitlenumber}{0}}%
                    \par\vspace{2em}
                    \phantomsection\addcontentsline{toc}{section}{\thetitlenumber\hspace{1em}\textbf{\@title}}%
                     }}
\makeatother



%% thm like environments
\let\question\relax
\let\endquestion\relax

\newtheoremstyle{QuestionStyle}{\topsep}{\topsep}%%% space between body and thm
		{}                      %%% Thm body font
		{}                              %%% Indent amount (empty = no indent)
		{\bfseries}            %%% Thm head font
		{)}                              %%% Punctuation after thm head
		{ }                           %%% Space after thm head
		{\thmnumber{#2}\thmnote{ \bfseries(#3)}}%%% Thm head spec
\theoremstyle{QuestionStyle}
\newtheorem{question}{}



\let\freeResponse\relax
\let\endfreeResponse\relax

%% \newtheoremstyle{ResponseStyle}{\topsep}{\topsep}%%% space between body and thm
%% 		{\wedn\bfseries}                      %%% Thm body font
%% 		{}                              %%% Indent amount (empty = no indent)
%% 		{\wedn\bfseries}            %%% Thm head font
%% 		{}                              %%% Punctuation after thm head
%% 		{3ex}                           %%% Space after thm head
%% 		{\underline{\underline{\thmname{#1}}}}%%% Thm head spec
%% \theoremstyle{ResponseStyle}

\usepackage[tikz]{mdframed}
\mdfdefinestyle{ResponseStyle}{leftmargin=1cm,linecolor=black,roundcorner=5pt,frametitlefont=\wedn\bfseries,%frametitle={\underline{\underline{Response}}:}
, font=\wedn\bfseries,}%\begin{mdframed}[style=mystyle]foo\end{mdframed}


\ifhandout
\NewEnviron{freeResponse}{}
\else
%\newtheorem{freeResponse}{Response:}
\newenvironment{freeResponse}{\begin{mdframed}[style=ResponseStyle]}{\end{mdframed}}
\fi



%% attempting to automate outcomes.

\newwrite\outcomefile
  \immediate\openout\outcomefile=\jobname.oc
\renewcommand{\outcome}[1]{\edef\theoutcomes{\theoutcomes #1~}%
\immediate\write\outcomefile{\unexpanded{\outcome}{#1}}}

%% \newcommand{\outcomelist}{\begin{itemize}\theoutcomes\end{itemize}}



%% my commands

\newcommand{\snap}{{\bfseries\itshape\textsf{Snap!}}}
\newcommand{\flavor}{\link[\snap]{https://snap.berkeley.edu/}}


\usepackage{tkz-euclide}
\tikzstyle geometryDiagrams=[rounded corners=.5pt,ultra thick,color=black]
\colorlet{penColor}{black} % Color of a curve in a plot


\author{Bart Snapp}

\title{Frieze patterns}

\begin{document}
\begin{abstract}
  We introduce \emph{frieze patterns}.
\end{abstract}
\maketitle

\begin{listOutcomes}
\item Describe what a frieze pattern is.
\item Identify minimal cells of frieze patterns.
\item Start to think about symmetry and frieze patterns.
\end{listOutcomes}

In design and architecture, a \index{frieze pattern}\emph{frieze},
pronounced ``freeze,'' is a horizontal band that is found near the
ceiling or roof. It is often decorated with a repeating design. Let's
see some examples:
\begin{center}
\includegraphics[width=.6\textwidth]{fpfrieze.pdf}
\end{center}
In each case, the pattern above can be thought of as an infinite strip
with the pattern continuing indefinitely.  Frieze patterns often have
symmetry in their structure. While in real-life, the symmetry would
not be perfect, it is useful in mathematics to simplify reality by
replacing the real patterns by idealized designs.  Today, we will
study the symmetry of frieze patterns.

\paragraph{Symmetries} Perhaps the most apparent symmetry that frieze
patters can exhibit is symmetry through \emph{horizontal
translations}:
\begin{center}
\includegraphics[width=.4\textwidth]{fptrans.pdf}
\end{center}
We will not only insist that a frieze pattern always have
translational symmetry, but we will also insist that such a
translation is \emph{discrete}, meaning that it occurs in distinct
``hops'' of some minimal distance. So, for example, the while the
following patterns have minimal discrete translations, 
\begin{center}
\includegraphics[width=.4\textwidth]{fptransDis.pdf}
\end{center}
this pattern
\begin{center}
\includegraphics[width=.4\textwidth]{fptransNot.pdf}
\end{center}
does not, and hence will not technically be considered a frieze
pattern. To be completely explicit: 
\begin{center}
\textbf{All frieze patterns have symmetry through a minimal discrete
  horizontal translation.}
\end{center}
\begin{definition}\index{minimal cell}
A \textbf{minimal cell} is the smallest slice of a frieze
pattern that will produce the entire frieze pattern through horizontal
translations.
\end{definition}

\begin{warning}
The minimal cell of a frieze pattern is not unique! Consider
\begin{center}
\includegraphics[width=.4\textwidth]{fptransMinEG.pdf}
\end{center}
both
\begin{center}
  \[
  \vcenter{\hbox{\includegraphics[width=.08\textwidth]{fptransMinEGCell.pdf}}} \qquad\text{and}\qquad  \vcenter{\hbox{\includegraphics[width=.08\textwidth]{fptransMinEGCell2.pdf}}}
  \]
\end{center}
are minimal cells.
\end{warning} 


\mynewpage



\begin{question} 
For each frieze pattern below, identify a MINIMAL CELL:
\[
\includegraphics[width=.6\textwidth]{fpfriezeHoriz.pdf}
\]
\begin{freeResponse}
    I've boxed the minimal cells below:
    \begin{center}
      \includegraphics[width=.6\textwidth]{fpfriezeHorizMINCELL.png}
    \end{center}
\end{freeResponse}
\end{question}
\mynewpage


\begin{question}
  Find THREE examples of frieze patterns on the INTERNET.
  Display them below.

  In EACH CASE, identify a minimal cell.
  \begin{freeResponse}
    \begin{enumerate}
    \item \begin{center}
      \includegraphics[width=.6\textwidth]{frieze1.png}
    \end{center}
    \item
      \begin{center}
      \includegraphics[width=.6\textwidth]{frieze2.png}
      \end{center}
    \item
      \begin{center}
        \includegraphics[width=.6\textwidth]{frieze3.png}
      \end{center}
    \end{enumerate}
  \end{freeResponse}
\end{question}
\mynewpage



\begin{question}
  LOOK at the frieze patterns above. DESCRIBE THREE different types of
  symmetry that some of these patterns have.
  \begin{freeResponse}
    Above I see that every frieze pattern has symmetry through
    translations.  Moreover, some frieze patterns have symmetry
    through
    \begin{enumerate}
    \item reflections over vertical lines,
    \item $180^\circ$ rotations, and
    \item reflections over horizontal
      lines.
    \end{enumerate}
  \end{freeResponse}
\end{question}


\end{document}
